\documentclass[10pt,a4paper]{amsart}
\usepackage{amsmath,amssymb,amsthm,graphicx}
\usepackage{verbatim}
\newenvironment{reference}{\comment}{\endcomment}
%\newenvironment{reference}{}{}
\newenvironment{slogan}{\comment}{\endcomment}
\newenvironment{history}{\comment}{\endcomment}

% For commutative diagrams we use Xy-pic
\usepackage[all]{xy}

% We use 2cell for 2-commutative diagrams.
\xyoption{2cell}
% Theorem environments.
%
\theoremstyle{plain}
\newtheorem{theorem}[subsection]{Theorem}
\newtheorem{proposition}[subsection]{Proposition}
\newtheorem{lemma}[subsection]{Lemma}

\theoremstyle{definition}
\newtheorem{definition}[subsection]{Definition}
\newtheorem{example}[subsection]{Example}
\newtheorem{exercise}[subsection]{Exercise}
\newtheorem{situation}[subsection]{Situation}

\theoremstyle{remark}
\newtheorem{remark}[subsection]{Remark}
\newtheorem{remarks}[subsection]{Remarks}

\numberwithin{equation}{subsection}

% Macros
%
\def\lim{\mathop{\mathrm{lim}}\nolimits}
\def\colim{\mathop{\mathrm{colim}}\nolimits}
\def\Spec{\mathop{\mathrm{Spec}}}
\def\Hom{\mathop{\mathrm{Hom}}\nolimits}
\def\Ext{\mathop{\mathrm{Ext}}\nolimits}
\def\SheafHom{\mathop{\mathcal{H}\!\mathit{om}}\nolimits}
\def\SheafExt{\mathop{\mathcal{E}\!\mathit{xt}}\nolimits}
\def\Sch{\mathit{Sch}}
\def\Mor{\mathop{\mathrm{Mor}}\nolimits}
\def\Ob{\mathop{\mathrm{Ob}}\nolimits}
\def\Sh{\mathop{\mathit{Sh}}\nolimits}
\def\NL{\mathop{N\!L}\nolimits}
\def\CH{\mathop{\mathrm{CH}}\nolimits}
\def\proetale{{pro\text{-}\acute{e}tale}}
\def\etale{{\acute{e}tale}}
\def\QCoh{\mathit{QCoh}}
\def\Ker{\mathop{\mathrm{Ker}}}
\def\Im{\mathop{\mathrm{Im}}}
\def\Coker{\mathop{\mathrm{Coker}}}
\def\Coim{\mathop{\mathrm{Coim}}}

% Boxtimes
%
\DeclareMathSymbol{\boxtimes}{\mathbin}{AMSa}{"02}

%
% Macros for moduli stacks/spaces
%
\def\QCohstack{\mathcal{QC}\!\mathit{oh}}
\def\Cohstack{\mathcal{C}\!\mathit{oh}}
\def\Spacesstack{\mathcal{S}\!\mathit{paces}}
\def\Quotfunctor{\mathrm{Quot}}
\def\Hilbfunctor{\mathrm{Hilb}}
\def\Curvesstack{\mathcal{C}\!\mathit{urves}}
\def\Polarizedstack{\mathcal{P}\!\mathit{olarized}}
\def\Complexesstack{\mathcal{C}\!\mathit{omplexes}}
% \Pic is the operator that assigns to X its picard group, usage \Pic(X)
% \Picardstack_{X/B} denotes the Picard stack of X over B
% \Picardfunctor_{X/B} denotes the Picard functor of X over B
\def\Pic{\mathop{\mathrm{Pic}}\nolimits}
\def\Picardstack{\mathcal{P}\!\mathit{ic}}
\def\Picardfunctor{\mathrm{Pic}}
\def\Deformationcategory{\mathcal{D}\!\mathit{ef}}
\usepackage{cleveref}
\crefname{originalsection}{original section}{original sections}
\Crefname{originalsection}{Original section}{Original sections}
\makeatletter
\newlabel{section-blowup-Z-first}{{0FBG}{source}{}{sourcepointer-1}{}}
\newlabel{section-blowup-Z-first@cref}{{[originalsection][0][]0FBG}{[1][0][]source}{}{}{}}
\makeatother
\crefname{originallemma}{original lemma}{original lemmas}
\Crefname{originallemma}{Original lemma}{Original lemmas}
\makeatletter
\newlabel{constructions-lemma-morphism-relative-proj}{{07ZG}{source}{}{sourcepointer-2}{}}
\newlabel{constructions-lemma-morphism-relative-proj@cref}{{[originallemma][0][]07ZG}{[1][0][]source}{}{}{}}
\makeatother
\crefname{originallemma}{original lemma}{original lemmas}
\Crefname{originallemma}{Original lemma}{Original lemmas}
\makeatletter
\newlabel{divisors-lemma-strict-transform}{{080E}{source}{}{sourcepointer-3}{}}
\newlabel{divisors-lemma-strict-transform@cref}{{[originallemma][0][]080E}{[1][0][]source}{}{}{}}
\makeatother
\crefname{originalreference}{original reference}{original references}
\Crefname{originalreference}{Original reference}{Original references}
\makeatletter
\newlabel{item-find-Z-in-blowup}{{0FE9}{source}{}{sourcepointer-4}{}}
\newlabel{item-find-Z-in-blowup@cref}{{[originalreference][0][]0FE9}{[1][0][]source}{}{}{}}
\makeatother
\crefname{originallemma}{original lemma}{original lemmas}
\Crefname{originallemma}{Original lemma}{Original lemmas}
\makeatletter
\newlabel{lemma-gysin-bivariant}{{0B79}{source}{}{sourcepointer-5}{}}
\newlabel{lemma-gysin-bivariant@cref}{{[originallemma][0][]0B79}{[1][0][]source}{}{}{}}
\makeatother
\crefname{originalreference}{original reference}{original references}
\Crefname{originalreference}{Original reference}{Original references}
\makeatletter
\newlabel{remark-restriction-bivariant}{{0F9Z}{source}{}{sourcepointer-6}{}}
\newlabel{remark-restriction-bivariant@cref}{{[originalreference][0][]0F9Z}{[1][0][]source}{}{}{}}
\makeatother
\crefname{originallemma}{original lemma}{original lemmas}
\Crefname{originallemma}{Original lemma}{Original lemmas}
\makeatletter
\newlabel{lemma-gysin-commutes-gysin}{{0B73}{source}{}{sourcepointer-7}{}}
\newlabel{lemma-gysin-commutes-gysin@cref}{{[originallemma][0][]0B73}{[1][0][]source}{}{}{}}
\makeatother
\title{The Gysin pullback of the normal cone of Z in X is the normal cone of Z in Y}
\author{The Stacks Project Authors}
\date{stacks-project:chow.tex:local-snapshot-20261012; GNU FDL 1.2 or later}
\begin{document}
\maketitle
\noindent\textit{Editorial note.} All delivered mathematical text is the author's own, copied from the pinned Stacks Project chapter file chow.tex; no statement or proof step was rewritten, reordered or removed.\par\medskip
\noindent\textit{Editorial note.} Preamble provenance: the pinned chapter file opens with an input of the shared preamble and carries no preamble text of its own, so the project's single Stacks-wide preamble staged at sources/stacks-chow/preamble.tex is excerpted verbatim as two spans, lines 14-18 (Xy-pic setup) and lines 154-217 (the native theorem declarations and mathematical macros). Nothing was added, removed or redefined; the wrapper class and the amsmath/amssymb/amsthm/graphicx base packages are supplied by the assembler.\par\medskip
\noindent\textit{Editorial note.} The counted statement and proof contain no cross reference, no citation and no included file, so the delivered body is self contained.\par\medskip
\noindent\textit{Editorial note.} Canonical location of the counted proposition: Stacks tag 0FEB (sources/stacks-chow/pinned-original-source.tex line 12260).\par\medskip
\noindent\textit{Allowed standard prerequisites (editorial).}\begin{itemize}
\item The blowing up of infinity(Z) and the rational map induced by the inclusion of ideals. Same chapter, Section labelled section-blowup-Z-first, invoked at proof lines 12268-12270; Constructions, Lemma labelled constructions-lemma-morphism-relative-proj, cited at proof lines 12276-12278. Produces the open U in W mapping to W'.
\item Strict transform of P\textasciicircum{}1\_Y inside W and its special fibre. Divisors, Lemma labelled divisors-lemma-strict-transform, cited at proof lines 12289-12292. Identifies W'' with the blowup of P\textasciicircum{}1\_Y in infinity(Z) and (W'' cap U)\_infinity with C\_Z Y.
\item Gysin maps of effective Cartier divisors and their commutation. Same chapter, Lemma labelled lemma-gysin-bivariant, cited at proof lines 12295-12296, and Lemma labelled lemma-gysin-commutes-gysin, cited at proof lines 12317-12319. Supplies the bivariant classes i\textasciicircum{}*, (i'\_infinity)\textasciicircum{}* and i\_infinity\textasciicircum{}* and lets their order be exchanged.
\item Vanishing at infinity and the factorization of the fibre of i over infinity. Same chapter, Remark labelled remark-restriction-bivariant, cited at proof lines 12298-12300, and the notation of the section at proof lines 12319-12321. Gives i\_infinity\textasciicircum{}* i\textasciicircum{}*[U] = [C\_Z Y]\_n and the description of the fibre through o.
\end{itemize}\medskip
\ifdefined\hypertarget\hypertarget{sourcepointer-1}{}\fi
\noindent\textit{Source reference.} Original section 0FBG: sources/stacks-chow/pinned-original-source.tex, source lines 11646--11646. Same-chapter pointer into the author environment lemma at source line 11646; invoked at proof line 12270 of the counted proof. It is retained only as a navigation pointer; every load-bearing step is inside the delivered span.\par\medskip
\ifdefined\hypertarget\hypertarget{sourcepointer-2}{}\fi
\noindent\textit{Source reference.} Original lemma 07ZG: sources/stacks-chow/tags.txt, source lines 10146--10146. Cross-chapter standard prerequisite located in the project tag list staged in the local snapshot; invoked at proof line 12279 of the counted proof. It is retained only as a navigation pointer; every load-bearing step is inside the delivered span.\par\medskip
\ifdefined\hypertarget\hypertarget{sourcepointer-3}{}\fi
\noindent\textit{Source reference.} Original lemma 080E: sources/stacks-chow/tags.txt, source lines 10179--10179. Cross-chapter standard prerequisite located in the project tag list staged in the local snapshot; invoked at proof line 12290 of the counted proof. It is retained only as a navigation pointer; every load-bearing step is inside the delivered span.\par\medskip
\ifdefined\hypertarget\hypertarget{sourcepointer-4}{}\fi
\noindent\textit{Source reference.} Original reference 0FE9: sources/stacks-chow/pinned-original-source.tex, source lines 11678--11678. Same-chapter pointer into the author environment lemma at source line 11678; invoked at proof line 12295 of the counted proof. It is retained only as a navigation pointer; every load-bearing step is inside the delivered span.\par\medskip
\ifdefined\hypertarget\hypertarget{sourcepointer-5}{}\fi
\noindent\textit{Source reference.} Original lemma 0B79: sources/stacks-chow/pinned-original-source.tex, source lines 6149--6149. Same-chapter pointer into the author environment lemma at source line 6149; invoked at proof line 12296 of the counted proof. It is retained only as a navigation pointer; every load-bearing step is inside the delivered span.\par\medskip
\ifdefined\hypertarget\hypertarget{sourcepointer-6}{}\fi
\noindent\textit{Source reference.} Original reference 0F9Z: sources/stacks-chow/pinned-original-source.tex, source lines 6184--6184. Same-chapter pointer into the author environment remark at source line 6184; invoked at proof line 12299 of the counted proof. It is retained only as a navigation pointer; every load-bearing step is inside the delivered span.\par\medskip
\ifdefined\hypertarget\hypertarget{sourcepointer-7}{}\fi
\noindent\textit{Source reference.} Original lemma 0B73: sources/stacks-chow/pinned-original-source.tex, source lines 5496--5496. Same-chapter pointer into the author environment lemma at source line 5496; invoked at proof line 12318 of the counted proof. It is retained only as a navigation pointer; every load-bearing step is inside the delivered span.\par\medskip
\begin{lemma}
\label{lemma-relation-normal-cones}
With notation as above we have
$$
o^*[C_ZX]_n = [C_Z Y]_{n - 1}
$$
in $\CH_{n - 1}(Y \times_{o, C_Y X} C_ZX)$.
\end{lemma}

\begin{proof}
Denote $W \to \mathbf{P}^1_X$ the blowing up of $\infty(Z)$ as in
Section \ref{section-blowup-Z-first}.
Similarly, denote $W' \to \mathbf{P}^1_X$ the blowing up of $\infty(Y)$.
Since $\infty(Z) \subset \infty(Y)$ we get an opposite inclusion
of ideal sheaves and hence a map of the graded algebras
defining these blowups. This produces a rational morphism from $W$
to $W'$ which in fact has a canonical representative
$$
W \supset U \longrightarrow W'
$$
See Constructions, Lemma \ref{constructions-lemma-morphism-relative-proj}.
A local calculation (omitted) shows that $U$ contains at least all points
of $W$ not lying over $\infty$ and the open subscheme $C_Z X$ of the special
fibre. After shrinking $U$ we may assume $U_\infty = C_Z X$ and
$\mathbf{A}^1_X \subset U$. Another local calculation (omitted)
shows that the morphism $U_\infty \to W'_\infty$
induces the canonical morphism $C_Z X \to C_Y X \subset W'_\infty$
of normal cones induced by the inclusion of ideals sheaves
coming from $Z \subset Y$. Denote $W'' \subset W$ the strict transform of
$\mathbf{P}^1_Y \subset \mathbf{P}^1_X$ in $W$. Then $W''$ is the blowing
up of $\mathbf{P}^1_Y$ in $\infty(Z)$ by
Divisors, Lemma \ref{divisors-lemma-strict-transform}
and hence $(W'' \cap U)_\infty = C_ZY$.

\medskip\noindent
Consider the effective Cartier divisor $i : \mathbf{P}^1_Y \to W'$
from (\ref{item-find-Z-in-blowup}) and its associated bivariant class
$i^* \in A^1(\mathbf{P}^1_Y \to W')$ from Lemma \ref{lemma-gysin-bivariant}.
We similarly denote $(i'_\infty)^* \in A^1(W'_\infty \to W')$ the
Gysin map at infinity. Observe that the restriction of $i'_\infty$
(Remark \ref{remark-restriction-bivariant}) to $U$ is the restriction of
$i_\infty^* \in A^1(W_\infty \to W)$ to $U$. On the one hand we have
$$
(i'_\infty)^* i^* [U]_{n + 1} =
i_\infty^* i^* [U]_{n + 1} =
i_\infty^* [(W'' \cap U)_\infty]_{n + 1} =
[C_ZY]_n
$$
because $i_\infty^*$ kills all classes supported over $\infty$, because
$i^*[U]$ and $[W'']$ agree as cycles over $\mathbf{A}^1$, and because
$C_ZY$ is the fibre of $W'' \cap U$ over $\infty$.
On the other hand, we have
$$
(i'_\infty)^* i^* [U]_{n + 1} =
i^* i_\infty^*[U]_{n + 1} =
i^* [U_\infty] =
o^*[C_YX]_n
$$
because $(i'_\infty)^*$ and $i^*$ commute
(Lemma \ref{lemma-gysin-commutes-gysin})
and because the fibre of $i : \mathbf{P}^1_Y \to W'$ over $\infty$
factors as $o : Y \to C_YX$ and the open immersion $C_YX \to W'_\infty$.
The lemma follows.
\end{proof}

\end{document}
